\subsection{Matrix Factorization and Topic Models}

Many unsupervised models express data as combinations of latent factors.

\begin{definition} \label{def:sl-low-rank-matrix-factorization} \hlt{Low-Rank Matrix Factorization}\\
Let $X\in\R^{m\times n}$ and $1\leq r\leq\min\{m,n\}$.\\
A \hlt{rank-$r$ matrix factorization} is pair $(U,V)\in\R^{m\times r}\times\R^{n\times r}$ with fitted matrix $\widehat{X}=UV^\top$ and $\operatorname{rank}(\widehat{X})\leq r$.\\
For a loss $\mathcal{L}$, choose $(U,V)$ to minimize $\mathcal{L}(X,UV^\top)$; the columns encode latent row and column factors.
\end{definition}

\begin{definition} \label{def:sl-nonnegative-matrix-factorization} \hlt{Nonnegative Matrix Factorization}\\
Let $X\in\R_+^{n\times p}$ and $1\leq r\leq\min\{n,p\}$.\\
A \hlt{nonnegative matrix factorization} is a pair $(W,H)\in\R_+^{n\times r}\times\R_+^{r\times p}$ with fitted matrix $\widehat{X}=WH$.\\
For loss $\mathcal{L}$, minimize $\mathcal{L}(X,WH)$ over $W,H\geq0$, yielding additive parts-based factors.
\end{definition}

\begin{method} \label{meth:sl-poisson-nmf-multiplicative-updates} \hlt{Poisson NMF Multiplicative Updates}\\
For count data $X=(x_{ij})$, set $\mu=WH$ and $L(W,H)=\sum\limits_{i,j}\{x_{ij}\log\mu_{ij}-\mu_{ij}\}$, up to constants.
\begin{enumerate}[label=(\roman*)]
\setlength{\itemsep}{0pt}
\item Initialize $W,H>0$.
\item Set $\rho^W_{ik}=(\sum\limits_j \frac{h_{kj}x_{ij}}{\mu_{ij}})/(\sum\limits_j h_{kj})$ and update $w_{ik}\gets w_{ik}\rho^W_{ik}$.
\item Recompute $\mu=WH$, set $\rho^H_{kj}=(\sum\limits_i \frac{w_{ik}x_{ij}}{\mu_{ij}})/(\sum\limits_i w_{ik})$, and update $h_{kj}\gets h_{kj}\rho^H_{kj}$.
\item Repeat (ii)--(iii) until $L(W,H)$ stabilises.
\end{enumerate}
The updates are local; different initialisations may give different NMF decompositions.
\end{method}

\begin{remark} \label{rem:sl-nmf-nonuniqueness} \hlt{NMF Non-Uniqueness}\\
An NMF $X=WH$ is not identifiable from nonnegativity alone: if $A$ is invertible and $WA\geq0$, $A^{-1}H\geq0$, then $X=(WA)(A^{-1}H)$. Interpretability therefore depends on additional geometry, constraints, or initialization.
\end{remark}

\begin{definition} \label{def:sl-archetypal-analysis} \hlt{Archetypal Analysis}\\
Let $X\in\R^{n\times p}$ have observed rows $x_i^T$, and fix $r$ archetypes.
\begin{enumerate}[label=(\roman*)]
\setlength{\itemsep}{0pt}
\item Observation weights $W=(w_{ik})\in\R_+^{n\times r}$ satisfy $\sum\limits_{k=1}^r w_{ik}=1$.
\item Archetype weights $B=(b_{ki})\in\R_+^{r\times n}$ satisfy $\sum\limits_{i=1}^n b_{ki}=1$.
\item Archetypes are the rows of $H=BX$, hence lie in $\operatorname{conv}\{x_1,\ldots,x_n\}$.
\item Archetypal analysis chooses $W,B$ minimising $J(W,B)=\norm{X-WBX}_F^2$.
\end{enumerate}
Archetypes are convex-hull extremes rather than central prototypes.
\end{definition}

\begin{remark} \label{rem:sl-archetypes-versus-kmeans-and-nmf} \hlt{Archetypes, $k$-Means, and NMF}\\
For prototype matrix $H$ and assignment/mixture weights $W$:
\begin{enumerate}[label=(\roman*)]
\setlength{\itemsep}{0pt}
\item $k$-means uses one-hot rows of $W$.
\item NMF uses nonnegative $W,H$ without requiring prototype rows of $H$ to be observed-data mixtures.
\item Archetypal analysis uses simplex rows of $W$ and constrains $H=BX$ with simplex rows of $B$.
\end{enumerate}
Archetypes are convex-hull prototypes; $k$-means is hard-assigned and NMF is hull-free.
\end{remark}

\begin{definition} \label{def:sl-topic-model} \hlt{Topic Model}\\
Fix $K$ topics and a vocabulary of size $p$.
\begin{enumerate}[label=(\roman*)]
\setlength{\itemsep}{0pt}
\item Topic $k$ is a word distribution $\beta_k\in\Delta_{p-1}$.
\item Document $d$ has topic mixture $\theta_d\in\Delta_{K-1}$.
\item Token $m$ in document $d$ is generated by $z_{dm}\sim\operatorname{Cat}(\theta_d)$ and $w_{dm}\sim\operatorname{Cat}(\beta_{z_{dm}})$.
\item Latent Dirichlet allocation adds priors $\theta_d\sim\operatorname{Dir}(\alpha)$ and $\beta_k\sim\operatorname{Dir}(\eta)$.
\end{enumerate}
\end{definition}

\begin{remark} \label{rem:sl-factorization-regularization} \hlt{Regularized Factors}\\
For a factorization $X\approx UV^\top$, constraints select an identifiable and interpretable solution class.
\begin{enumerate}[label=(\roman*)]
\setlength{\itemsep}{0pt}
\item Low rank bounds latent dimension.
\item Sparsity restricts active variables or factors.
\item Nonnegativity enforces additive parts.
\item Probabilistic priors encode uncertainty and shrinkage.
\end{enumerate}
These restrictions convert an underdetermined algebraic fit into a statistical learning problem.
\end{remark}
